% GENERATED, edit the French unit or the translation table
\mxunit{premiere-c02}{Quadratics: forms and choice of method}\label{C02}
\begin{keypoint}[Chapter path]
\textbf{Objectives.} Choose a form; solve an equation or an inequality; read an extremum; prove the formula for the roots.

\textbf{Prerequisites.} Special products, zero-product property, sign of a product, square roots.

\textbf{Yearly reference point.} 22 September--1 October: 8 hours\end{keypoint}
\section*{Diagnostic: six questions, 10 minutes}\label{C02-D}
Answer without consulting the lesson. Mark the questions for which help is needed. This diagnostic is not graded.
\begin{enumerate}
\item Expand $(x-3)^2$.
\item Factorise $x^2-16$.
\item Solve $(x-2)(x+5)=0$.
\item Solve $x^2=9$ in $\R$.
\item Give the sign of $2x-6$.
\item Find the minimum of $(x-1)^2+4$.
\end{enumerate}
\paragraph{After the diagnostic.} Revisit the flagged topics with the teacher; redo a similar question, then explain the rule used.
\paragraph{Speaking and logic.} Present two methods for factorising a quadratic. In the proof using the discriminant, name the three cases and explain why they are exhaustive.
\mxnewpage
\section*{Investigate before formalising}\label{C02-RECH}
A rectangle has an area of 24 m$^2$ and a perimeter of 20 m. Without trying every possible dimension at random, determine its sides. Start with a drawing and explain the two constraints.
\paragraph{Organisation.} Individual investigation 5 min, comparison in pairs 5 min, class discussion 5 min. Keep a record of one attempt, even an unfinished one. A conjecture must then be justified.
\mxnewpage
\section{Three forms, three readings}\label{C02-U1}
\begin{keypoint}[Definition and results]
A quadratic function is written $f(x)=ax^2+bx+c$ with $a\ne0$. The expanded form makes it easy to calculate $f(0)$; the vertex form $a(x-\alpha)^2+\beta$ reveals the vertex and the extremum; a factorised form $a(x-r_1)(x-r_2)$ reveals the roots and the sign. The roots are the solutions of $f(x)=0$.
\end{keypoint}
\paragraph{Status and justification.} Completing the square relies on $x^2+2ux=(x+u)^2-u^2$. The existence of the vertex form is proved by expansion; its general computation is not an expected routine.
\begin{example}$f(x)=x^2-6x+5=(x-3)^2-4=(x-1)(x-5)$. The minimum is $-4$ at 3; the roots are 1 and 5; the y-intercept is 5.\end{example}\label{C02-EX1}
\mxinclude{figures/C02-U1}
\begin{pitfall}A factorised expression is not obtained by putting each term separately in brackets.\end{pitfall}
\begin{application}Give three forms of $g(x)=x^2+2x-3$. Read off its roots and its minimum.\end{application}\label{C02-AP1}
\section{Discriminant and solving}\label{C02-U2}
\begin{keypoint}[Definition and results]
For $ax^2+bx+c=0$ with $a\ne0$, let $\Delta=b^2-4ac$. If $\Delta<0$, there is no real root; if $\Delta=0$, there is one root $-\dfrac{b}{2a}$; if $\Delta>0$, there are two roots $\dfrac{-b-\sqrt\Delta}{2a}$ and $\dfrac{-b+\sqrt\Delta}{2a}$. To order the roots, take the sign of a into account.
\end{keypoint}
\paragraph{Status and justification.} Multiplying the equation by $4a\ne0$ gives $(2ax+b)^2=\Delta$. A square is positive or zero: the three cases follow. If $\Delta>0$, $2ax+b=\pm\sqrt\Delta$. Splitting into cases avoids taking the square root of a negative number.
\begin{example}$2x^2+x-3=0$: $\Delta=1+24=25$; $x=\dfrac{-1\pm5}{4}$, that is $1$ and $-\dfrac{3}{2}$. Check: $(x-1)(2x+3)=2x^2+x-3$.\end{example}\label{C02-EX2}
\begin{pitfall}The denominator is $2a$, not 2; the sign of b is kept with its brackets in $b^2$.\end{pitfall}
\begin{application}Solve $3x^2-6x+3=0$ and $x^2+2x+5=0$.\end{application}\label{C02-AP2}
\section{Factorising without automating the discriminant}\label{C02-U3}
\begin{keypoint}[Definition and results]
If two real roots exist, their sum is $-\dfrac{b}{a}$ and their product is $\dfrac{c}{a}$. An obvious root or a special product often allows a shorter factorisation. The polynomials of degree 2 that vanish at two distinct reals r and s are exactly $a(x-r)(x-s)$ with $a\ne0$.
\end{keypoint}
\paragraph{Status and justification.} Expanding $a(x-r)(x-s)=ax^2-a(r+s)x+ars$, identifying the coefficients gives the relations. Conversely, if r and s have this sum and this product, the expansion recovers the polynomial.
\begin{example}$x^2-7x+12$: 3 and 4 have sum 7 and product 12. Hence $(x-3)(x-4)$. To impose in addition $f(0)=24$, we need $12a=24$, so $a=2$.\end{example}\label{C02-EX3}
\begin{pitfall}The sum of the roots is $-\dfrac{b}{a}$; do not forget the minus sign. A single root is not enough to determine a.\end{pitfall}
\begin{application}Determine the polynomial f of degree 2 such that $f(1)=f(4)=0$ and $f(0)=-8$.\end{application}\label{C02-AP3}
\section{Sign, inequalities and optimisation}\label{C02-U4}
\begin{keypoint}[Definition and results]
For two distinct roots $r<s$, the quadratic has the sign of a outside $\intff{r}{s}$ and the opposite sign between the roots. It is zero at the roots. With no real root, it keeps the sign of a. With a double root, it keeps the sign of a away from that root. Optimisation is done on the domain of the problem, which may limit access to the vertex.
\end{keypoint}
\paragraph{Status and justification.} The sign table of the factors of $a(x-r)(x-s)$ gives the result. The vertex form compares $f(x)$ with $\beta$ through the sign of $a(x-\alpha)^2$.
\begin{example}$-2(x-1)(x-5)\geq0$ is equivalent to $x\in\intff{1}{5}$. For $A(x)=x(12-x)$ on $\intff{0}{12}$, $A(x)=36-(x-6)^2\leq36$, with equality at 6. A rectangle with perimeter 24 m therefore has a maximum area of 36 m$^2$ when it is a square.\end{example}\label{C02-EX4}
\begin{pitfall}In a strict inequality, the roots are excluded. Check the geometric bounds before concluding.\end{pitfall}
\begin{application}Solve $x^2-4x-5<0$, then maximise $x(8-x)$ on $\intff{0}{3}$.\end{application}\label{C02-AP4}
\mxnewpage
\section*{Exercise path}\label{C02-PARC}
Start with exercises 1 and 2. Then choose one of exercises 3 or 4 and do it without help. Exercises 5 to 7 are a reserve for transfer, investigation or programming: the teacher designates one. Exercise 8 is revisited after a delay.
\begingroup
\mxexercisehead{C02-E1}{Completing the square}
Consider the function $f$ defined on $\R$ by $f(x)=x^2-8x+7$.
\begin{enumerate}[label=\arabic*.]
\item Complete the square to write $f(x)$ in the form $(x-4)^2+k$, where $k$ is a real number to be determined.
\item Deduce a factorisation of $f(x)$ using a special product.
\item Draw the sign table of $f$, then solve the inequality $f(x)\leq0$ in $\R$.
\end{enumerate}
\mxexercisehead{C02-E2}{Discussing without heavy calculation}
We study the functions $f(x)=x^2+4$, $g(x)=(x+2)^2$ and $h(x)=-3(x-1)(x+4)$, all defined on $\R$.
\begin{enumerate}[label=\arabic*.]
\item For each function, determine any real roots by using the given form directly.
\item State the sign of each function on $\R$. State explicitly the values where it vanishes.
\item Explain why calculating a discriminant is necessary in none of the three cases.
\end{enumerate}
\mxexercisehead{C02-E3}{Choosing a strategy}
We wish to solve the following three equations in $\R$:\[x^2-9=0,\qquad x^2-9x+20=0,\qquad 2x^2-3x-1=0.\]
\begin{enumerate}[label=\arabic*.]
\item For each equation, choose a method: special product, finding two numbers with a known sum and product, or discriminant. Explain this choice.
\item Solve the three equations and give the solution sets in exact form.
\item Check the solutions of the first two equations by substitution. For the third, check the sum and the product of the roots obtained.
\end{enumerate}
\mxexercisehead{C02-E4}{Recovering a parabola}
The curve of a polynomial function $f$ of degree 2 crosses the x-axis at the points with x-coordinates $-2$ and 3. It also passes through the point $A\pt{1;12}$.
\begin{enumerate}[label=\arabic*.]
\item Determine an expression for $f(x)$ from these three pieces of information. State the form chosen and why it is suitable.
\item Expand the expression found. Deduce the y-intercept of the curve.
\item Check that the function obtained satisfies the three initial pieces of information.
\end{enumerate}
\mxexercisehead{C02-E5}{A vegetable garden against a wall}
A gardener has 20 m of fencing to enclose three sides of a rectangular vegetable garden. The fourth side runs along a wall and needs no fencing. Let $x$ be the length, in metres, of each of the two sides perpendicular to the wall. All the fencing is used.
\begin{enumerate}[label=\arabic*.]
\item Make a sketch and express the length of the side parallel to the wall in terms of $x$. State the values of $x$ for which the rectangle is not degenerate.
\item Express the area $A(x)$, in square metres. Determine its maximum value and the corresponding dimensions, using the method of your choice.
\item Check that these dimensions use exactly 20 m of fencing.
\end{enumerate}
\mxinclude{figures/C02-E5}
\mxexercisehead{C02-E6}{A parameter}
For each real number $m$, consider the equation in the real unknown $x$:\[x^2-2x+m=0.\]The number $m$ is a parameter: it stays fixed while the equation is solved.
\begin{enumerate}[label=\arabic*.]
\item Determine the values of $m$ for which the equation has two distinct real solutions.
\item Determine the values of $m$ for which it has a single real solution, then those for which it has none.
\item Give the solution in the case where it is unique. Explain why the three cases considered cover all values of $m$.
\end{enumerate}
\mxexercisehead{C02-E7}{Link with sequences}
A sequence is defined by $u_n=n^2-5n+4$ for every natural number $n$. We look for all the indices at which the term of the sequence is strictly negative.
\begin{enumerate}[label=\arabic*.]
\item First solve the inequality $x^2-5x+4<0$ in $\R$.
\item Deduce all the natural numbers $n$ such that $u_n<0$. Justify why a real interval is not the final answer here.
\end{enumerate}
\mxexercisehead{C02-E8}{Delayed review}
A student studies the function $f$ defined on $\R$ by $f(x)=2(x+3)^2-5$. She writes: ``The minimum is $-5$ and it is reached at $x=3$.''
\begin{enumerate}[label=\arabic*.]
\item State what is correct and what is wrong in this sentence. Rewrite it correctly, justifying your answer.
\item Solve the inequality $f(x)\leq-3$ in $\R$, then check the two endpoints of the set obtained.
\end{enumerate}
\paragraph{Homework.} Do only the task designated at the end of the session, in the dated progression: 10--20 min at most. The exercises depend on the topics already studied, so some questions are deferred. Revisit exercise 8 one week after the chapter (10 min). If the allotted time is exceeded, note the difficulty and bring your attempts; do not pile up the whole reserve.
\endgroup
\mxneed{12}\section*{To do in the CQFD textbook}
These exercises are done in your textbook. The progression indicates the session; they replace the homework planned for that day (10 to 20 minutes in total).
\paragraph{CQFD, exercise 43 p.~99, questions a and c.}\label{C02-M1} First find an obvious root, then the other using the product of the roots.
\paragraph{CQFD, exercise 49 p.~99, questions a and b.}\label{C02-M2} Translate the sum and the product into an equation; do not assume in advance that the numbers exist.
\begingroup\small
\mxnewpage\section*{MCQ 1: understand and choose}\label{C02-Q1}
Five single-answer questions for practice. Choose one answer per question; then explain one distractor. This chapter format does not reproduce the official EAM format.
\begin{enumerate}
\item The roots of $(x+1)(x-4)$ are
\begin{enumerate}[label=\Alph*.,itemsep=2pt]
\item $1$ and $4$
\item $-1$ and $-4$
\item $-1$ and $4$
\item $1$ and $-4$
\end{enumerate}
\item The discriminant of $x^2-2x+3$ equals
\begin{enumerate}[label=\Alph*.,itemsep=2pt]
\item $-8$
\item $8$
\item $16$
\item $-2$
\end{enumerate}
\item The vertex of $(x-2)^2+5$ is
\begin{enumerate}[label=\Alph*.,itemsep=2pt]
\item $\pt{2;-5}$
\item $\pt{-2;5}$
\item $\pt{5;2}$
\item $\pt{2;5}$
\end{enumerate}
\item The sum of the roots of $2x^2-6x+1$ equals
\begin{enumerate}[label=\Alph*.,itemsep=2pt]
\item $6$
\item $3$
\item $-3$
\item $\dfrac{1}{2}$
\end{enumerate}
\item $x^2-4<0$ is equivalent to
\begin{enumerate}[label=\Alph*.,itemsep=2pt]
\item $x>2$
\item $x<-2$
\item $x\in\intoo{-2}{2}$
\item $x\in\intff{-2}{2}$
\end{enumerate}
\end{enumerate}
\endgroup
\begingroup\small
\mxneed{18}\section*{MCQ 2: understand and choose}\label{C02-Q2}
Five single-answer questions for practice. Choose one answer per question; then explain one distractor. This chapter format does not reproduce the official EAM format.
\begin{enumerate}
\item $x^2+6x+9$ has as its factorisation
\begin{enumerate}[label=\Alph*.,itemsep=2pt]
\item $(x-3)^2$
\item $(x+9)(x+1)$
\item $(x+6)^2$
\item $(x+3)^2$
\end{enumerate}
\item The product of the roots of $3x^2+x-2$ equals
\begin{enumerate}[label=\Alph*.,itemsep=2pt]
\item $-\dfrac{2}{3}$
\item $-\dfrac{1}{3}$
\item $\dfrac{2}{3}$
\item $-2$
\end{enumerate}
\item $-x^2+4x$ reaches its maximum on $\intff{0}{1}$ at
\begin{enumerate}[label=\Alph*.,itemsep=2pt]
\item $2$
\item $0$
\item $1$
\item $4$
\end{enumerate}
\item A double root requires
\begin{enumerate}[label=\Alph*.,itemsep=2pt]
\item $a=0$
\item $\Delta=0$
\item $b=0$
\item $c=0$
\end{enumerate}
\item The polynomials of degree 2 with roots 0 and 2 are
\begin{enumerate}[label=\Alph*.,itemsep=2pt]
\item $a(x+2)$
\item $a(x-2)^2$
\item $x(x-2)$ only
\item $ax(x-2)$, $a\ne0$
\end{enumerate}
\end{enumerate}
\subsection*{True or false? Justify.}\label{C02-VF}
\begin{enumerate}
\item A quadratic function always has two real roots.
\item A quadratic with two distinct roots is negative between them.
\item For every real number x, $(x-2)^2\geq0$.
\end{enumerate}
\begin{keypoint}[Self-assessment]I justify the choice of form; I distinguish roots from vertex; I finish an inequality with a set; I respect the domain of the model.

For each objective, indicate: achieved without help; achieved with help; to be reworked. Write one example and one mistake to avoid.\end{keypoint}
\endgroup
\ifmxteacher\mxinclude{corriges/C02}\fi
